\documentclass[11pt,a4paper]{article}
\usepackage[utf8]{inputenc}
\usepackage[vietnamese]{babel}
\usepackage{amsmath,amssymb}
\usepackage{geometry}
\geometry{a4paper, total={180mm,260mm}, left=15mm, right=15mm, top=18mm, bottom=18mm}

\usepackage{tikz}
\usetikzlibrary{shapes.geometric, arrows.meta, positioning, calc}
\tikzset{
  arr/.style={-{Stealth[scale=0.8]}},
  every node/.append style={align=center}
}
\usepackage{tcolorbox}
\tcbuselibrary{skins,breakable}
\usepackage{fontawesome}
\usepackage{fancyhdr}
\usepackage{enumitem}

\pagestyle{fancy}
\fancyhf{}
\renewcommand{\headrulewidth}{0.5pt}
\renewcommand{\footrulewidth}{0.5pt}
\fancyhead[L]{\small\bfseries\color{violet!85!black}\faBook\ SỔ TAY KIẾN THỨC TOÁN 2}
\fancyhead[R]{\small\bfseries\color{violet!85!black}Lớp Toán Thầy Lĩnh -- loptoanthaylinh.vn}
\fancyfoot[C]{\small\color{gray}-- Sổ tay kiến thức: Bài 38. Thừa số, tích --}

\setlength{\parindent}{0pt}
\setlength{\parskip}{0pt}

\begin{document}

% --- BANNER BÀI HỌC ---
\begin{tcolorbox}[
  enhanced,
  colback=violet!8!white,
  colframe=violet!70!black,
  arc=3mm, boxrule=1.2pt,
  left=3mm, right=3mm, top=2mm, bottom=2mm,
  halign=center
]
  {\LARGE\bfseries\color{violet!90!black}\faBookmark\ BÀI 38: THỪA SỐ, TÍCH}\\[1mm]
  {\itshape\footnotesize\color{black!75} Tên gọi các thành phần của phép nhân và quy tắc đổi chỗ thừa số trên 1 trang A4}
\end{tcolorbox}

\vspace{-2mm}

\noindent
\begin{minipage}[t]{0.49\textwidth}\vspace{0pt}
\begin{tcolorbox}[
  enhanced,
  colback=violet!4!white,
  colframe=violet!75!black,
  arc=2mm, boxrule=0.8pt,
  left=2mm, right=2mm, top=1.8mm, bottom=1.8mm,
  title={\footnotesize\bfseries\color{white}\faKey\ 1. Thành phần phép nhân}
]
{\footnotesize
\begin{itemize}\setlength{\itemsep}{1pt}\setlength{\leftskip}{-2.5mm}
  \item \textbf{Cấu trúc tổng quát:}
  $$\mathbf{a \times b = c}$$
  \begin{itemize}\setlength{\itemsep}{0.5pt}\setlength{\leftskip}{-2.5mm}
    \item $a$: \textbf{Thừa số}.
    \item $b$: \textbf{Thừa số}.
    \item $c$: \textbf{Tích} (kết quả phép nhân).
  \end{itemize}
  \item \textbf{Lưu ý đặc biệt:}\\
  Biểu thức $a \times b$ cũng được gọi là một \textbf{"tích"}.\\
  \textit{Ví dụ:} Tích của $2$ và $5$ là $2 \times 5$ hay $10$.
\end{itemize}
}
\end{tcolorbox}
\end{minipage}
\hfill
\begin{minipage}[t]{0.49\textwidth}\vspace{0pt}
\begin{tcolorbox}[
  enhanced,
  colback=teal!4!white,
  colframe=teal!75!black,
  arc=2mm, boxrule=0.8pt,
  left=2mm, right=2mm, top=1.8mm, bottom=1.8mm,
  title={\footnotesize\bfseries\color{white}\faSliders\ 2. Tính chất giao hoán}
]
{\footnotesize
\begin{itemize}\setlength{\itemsep}{1pt}\setlength{\leftskip}{-2.5mm}
  \item Khi đổi chỗ các thừa số trong một tích thì \textbf{tích không thay đổi}:
  $$\mathbf{a \times b = b \times a}$$
  \item \textbf{Ví dụ minh họa:}\\
  $2 \times 5 = 10 \quad\text{và}\quad 5 \times 2 = 10$.\\
  $\implies 2 \times 5 = 5 \times 2$.
  \item Ý nghĩa: Giúp học sinh học thuộc nhanh bảng cửu chương (học $2 \times 5$ biết ngay $5 \times 2$).
\end{itemize}
}
\end{tcolorbox}
\end{minipage}

\vspace{-1.5mm}
\begin{tcolorbox}[
  enhanced,
  colback=orange!4!white,
  colframe=orange!80!black,
  arc=2mm, boxrule=0.8pt,
  left=2mm, right=2mm, top=1.5mm, bottom=1.5mm,
  title={\footnotesize\bfseries\color{white}\faSitemap\ 3. Sơ đồ tư duy định danh các thành phần phép nhân}
]
\centering
\begin{tikzpicture}[x=1.5cm, y=0.8cm]
  \node[rectangle, rounded corners=2pt, draw=violet!80!black, fill=violet!15, font=\scriptsize\bfseries] (ts1) at (0,0) {THỪA SỐ\\($3$)};
  \node[font=\normalsize\bfseries] at (1,0) {$\times$};
  \node[rectangle, rounded corners=2pt, draw=violet!80!black, fill=violet!15, font=\scriptsize\bfseries] (ts2) at (2,0) {THỪA SỐ\\($4$)};
  \node[font=\normalsize\bfseries] at (3,0) {$=$};
  \node[rectangle, rounded corners=2pt, draw=green!80!black, fill=green!15, font=\scriptsize\bfseries] (tich) at (4.5,0) {TÍCH\\($12$)};

  \draw[<->, thick, orange!80!black] (ts1) to[bend left=30] node[above, font=\tiny\bfseries] {Đổi vị trí: $3 \times 4 = 4 \times 3$} (ts2);
\end{tikzpicture}
\end{tcolorbox}

\vspace{-1.5mm}
\noindent
\begin{minipage}[t]{0.49\textwidth}\vspace{0pt}
\begin{tcolorbox}[
  enhanced,
  colback=green!4!white,
  colframe=green!60!black,
  arc=2mm, boxrule=0.8pt,
  left=2mm, right=2mm, top=1.8mm, bottom=1.8mm,
  title={\footnotesize\bfseries\color{white}\faCheckCircle\ 4. Ví dụ thực hành chuẩn mực}
]
{\footnotesize
\textbf{1. Gọi tên thành phần:}\\
Trong phép tính $5 \times 4 = 20$:\\
- $5$ và $4$ là các \textbf{thừa số}.\\
- $20$ là \textbf{tích} (và $5 \times 4$ cũng là một tích).
\textbf{2. Tính tích khi biết hai thừa số:}\\
Biết thừa số thứ nhất là $2$, thừa số thứ hai là $7$. Tính tích?\\
\textit{Giải:} Tích là: $2 \times 7 = \mathbf{14}$.
}
\end{tcolorbox}
\end{minipage}
\hfill
\begin{minipage}[t]{0.49\textwidth}\vspace{0pt}
\begin{tcolorbox}[
  enhanced,
  colback=red!4!white,
  colframe=red!75!black,
  arc=2mm, boxrule=0.8pt,
  left=2mm, right=2mm, top=1.8mm, bottom=1.8mm,
  title={\footnotesize\bfseries\color{white}\faExclamationTriangle\ 5. Mẹo tránh bẫy trắc nghiệm}
]
{\footnotesize
\begin{itemize}\setlength{\itemsep}{1pt}\setlength{\leftskip}{-2.5mm}
  \item \textbf{Bẫy yêu cầu "Tính tích":}\\
  Đề bài: "Tính tích của 4 và 3" $\implies$ Làm tính \textbf{nhân} ($4 \times 3 = 12$). Không được làm cộng hay trừ!
  \item Trong phép nhân, cả hai số đứng trước và sau dấu $\times$ đều có tên chung là \textbf{Thừa số}.
\end{itemize}
}
\end{tcolorbox}
\end{minipage}


\end{document}
